% =========================================================
% config/packages.tex
% SELURUH package buku didefinisikan di sini.
% Bab dan section TIDAK BOLEH memuat \usepackage.
% =========================================================

% ---------------------------------------------------------
% Bahasa dan encoding
% ---------------------------------------------------------
\usepackage[T1]{fontenc}
\usepackage[utf8]{inputenc}
\usepackage{lmodern}          % font Latin Modern (scalable, wajib untuk microtype)
\usepackage[indonesian]{babel}

% ---------------------------------------------------------
% Tata letak halaman
% ---------------------------------------------------------
\usepackage[
    a4paper,
    top=3cm,
    bottom=2.8cm,
    left=3cm,
    right=2.6cm,
    headsep=0.6cm,
    footskip=1.0cm
]{geometry}

% ---------------------------------------------------------
% Matematika
% ---------------------------------------------------------
\usepackage{amsmath}
\usepackage{amssymb}
\usepackage{amsthm}
\usepackage{mathtools}

% ---------------------------------------------------------
% Grafik dan diagram
% ---------------------------------------------------------
\usepackage{graphicx}
\usepackage{float}
\usepackage{tikz}
\usetikzlibrary{
    arrows.meta,
    positioning,
    shapes.geometric,
    shapes.misc,
    calc,
    fit,
    backgrounds,
    patterns,
    decorations.pathreplacing,
    matrix,
    chains
}

% ---------------------------------------------------------
% Tabel
% ---------------------------------------------------------
% Catatan: colortbl TIDAK dipakai. Pada instalasi MiKTeX tertentu
% colortbl versi baru memerlukan makro \insert@pcolumn yang belum
% tersedia pada array.sty bawaan (tools), sehingga tabel gagal
% dirakit. Pewarnaan baris tabel dihindari; gunakan booktabs dan
% judul kolom tebal.
\usepackage{booktabs}
\usepackage{longtable}
\usepackage{array}
\usepackage{tabularx}
\usepackage{multirow}
\usepackage{makecell}

% ---------------------------------------------------------
% Warna dan kotak
% ---------------------------------------------------------
\usepackage{xcolor}
\usepackage[most]{tcolorbox}

% ---------------------------------------------------------
% Daftar
% ---------------------------------------------------------
\usepackage{enumitem}

% ---------------------------------------------------------
% Kode program
% ---------------------------------------------------------
\usepackage{listings}

% ---------------------------------------------------------
% Keterangan (caption) dan judul bab/section
% ---------------------------------------------------------
\usepackage{caption}
\usepackage{titlesec}

% ---------------------------------------------------------
% Header dan footer
% ---------------------------------------------------------
\usepackage{fancyhdr}

% ---------------------------------------------------------
% Spasi
% ---------------------------------------------------------
\usepackage{setspace}
\usepackage{microtype}

% ---------------------------------------------------------
% Simbol
% ---------------------------------------------------------
\usepackage{pifont}

% ---------------------------------------------------------
% Referensi dan tautan (dimuat menjelang akhir)
% ---------------------------------------------------------
\usepackage[hidelinks,breaklinks]{hyperref}
\usepackage{url}

% Sitasi memakai gaya numerik [1], [2], ... mengikuti penomoran
% pada Daftar Referensi RPS. Gaya author-year tidak dipakai karena
% sebagian sumber resmi (misalnya manual laboratorium dan
% dokumentasi perangkat lunak) tidak mencantumkan tahun terbit,
% sehingga label author-year menjadi tidak lengkap.
\usepackage[numbers,sort&compress]{natbib}

% ---------------------------------------------------------
% Kompilasi bab mandiri
% ---------------------------------------------------------
\usepackage{subfiles}
